\section{Suma del resto en la llegada}\label{app:arrival}
Un solo cuadrado puede contribuir a la escala \(1/n\) que interesa en \eqref{eq:signed-mean}.
Al sustraer \(\Gamma\), cada soporte que sobrevive debe satisfacer las restricciones geométricas
de dos cuadrados. Ese coste adicional permite sumar los errores con una cota que decrece con \(b\).
Veámoslo primero en la llegada histórica corta \(H_j\), bajo \(\nu\). Para un conjunto
de casillas \(Y\), activamos sólo las restricciones diagonales soportadas en \(Y\), manteniendo
la máscara de torres completa. Llamemos \(e(Y)\) al logaritmo de la probabilidad de evitación
menos el cambio de predictor, y \(e_C(Y)\) a la misma expresión usando sólo los pares de \(C\).
El coeficiente de un soporte \(U\) es \(\sum_{Y\subseteq U}(-1)^{|U|-|Y|}e(Y)\).
Una casilla aislada se cancela al emparejar subconjuntos que difieren sólo en él. Si todos los
pares pertenecen a un cuadrado, el coeficiente se cancela con el de \(e_C\).
Si intervienen dos cuadrados distintos, ningún \(e_C\) contribuye: las casillas tendrían que
estar en los cuatro vértices de \(C\), cuyas únicas parejas diagonales son las dos de \(C\).
El resto suma por tanto los coeficientes originales que carecen de casillas aisladas y contienen
al menos dos cuadrados distintos. Lo sustraído tiene grado a lo sumo cuatro.

Sean \(v\) el número de casillas, \(h\) el total de etiquetas distintas de filas y columnas,
\(D\) el número de coordenadas libres y \(p\) el exponente del soporte. Prescribir las
\(h\) líneas da el factor \((j/n)^h\), junto con \(\delta^{-2v}\exp(C_Lv^2/j)\) en la
cota de prescripción. Para cada tipo, el número de posiciones es de orden \(n^D\).
Escribamos \(d_*=h-D\), \(q=p-D-1\). Dos cuadrados distintos imponen dos ecuaciones
lineales no proporcionales sobre las coordenadas de fila y columna. Así \(d_*\ge2\);
la cota del soporte da \(q\ge0\). El factor escalado es
\[
 n\,n^D j^{-p}(j/n)^h=(j/n)^{d_*-1}j^{-q}
 \le (b/\delta_*)\delta_*^{d_*-1}j^{-q}
 \qquad(0<b\le\delta_*).
\]
Extraemos \(b/\delta_*\) antes de sumar. Para \(v\le4\) el número de tipos tiene una cota
absoluta. Para los tipos no colineales el recuento ponderado usa \(d_*-1+q\ge(v-4)/3\).
De ahí que para \(5\le v\le N=\lceil4\log n/\log\rho\rceil\), fijado \(\rho\ge2\), la cota
tras sumar los tipos de ese tamaño sea
\[
 (b/\delta_*) C_L(C_1\delta^{-2})^v v^{d_{\rm pol}}\delta_*^{(v-4)/3}.
\]
Se elige \(\delta_*\) de manera que \(C_1\delta^{-2}\delta_*^{1/3}<1\).
En un tipo colineal \(v\ge3\) se tiene \(p=v+2,h=2v\), con a lo sumo \(4n\,n^v/v!\)
soportes. Resulta \(C(C_2\delta^{-2})^v v^2(j/n)^{v-2}\); se reduce \(\delta_*\)
una vez para sumar también esta cota. La sustitución de átomos hasta grado \(N\) cuesta
\(O(j^{-8})\) en valor, y la norma global acota la cola por \(O_\rho(n^{-2})\).
Ambos errores desaparecen tras escalar para cada \(b\) fijo. Así
\(\limsup_n n\mathbb E_\nu[|g^\flat|;H_j]\le C_\delta b\), con \(C_\delta\)
independiente de \(b\).

La cota obtenida controla una media de \(|g^\flat|\) bajo \(\nu\). En \eqref{eq:signed-history}
interviene además la exponencial bajo \(\mu\). El factor \(X\) de la densidad permite estimar
ambas con la misma cota, pues en la llegada se tiene
\[
 X\,|e^{-g^\flat}-1|=e^{g_T+\Gamma}|1-e^{g^\flat}|.
\]
En \(H_j\) la parte positiva del error diagonal tiene una cota de clase,
\(|g_T|+|\Gamma|\le C/j\) y \(|A-R_Q(n)|=O_{b,\delta}(1/n)\).
Se aplica \(|e^y-1|\le e^{y^+}|y|\) a la densidad \(Xe^{A-R_Q(n)}\).
En \(F_b\setminus H_j\) la probabilidad ponderada es \(O(e^{-2j})\).
Para \(0<X\le1\) vale \(X|\log X|\le1/e\); por encima de uno se usa
\(X|\log X|\le X(C_0+\beta j)\) en la clase exterior. El aporte de este suceso al lado
izquierdo de \eqref{eq:arrival-estimate} antes de escalar es \(O((1+j)e^{-2j})\).
Para obtener \eqref{eq:arrival-estimate} fijamos primero \(\delta\), los márgenes, \(\rho\) y
las constantes de clase, después \(\delta_*\) y por último el corte. El umbral de \(n\) puede
depender de \(b\). La constante \(C_\delta\) sirve para todos esos cortes, como exige la
comparación final de dos cortes en el apéndice~\ref{app:closing}.
